\documentclass[10pt,a4paper]{amsart}
\usepackage[margin=19mm]{geometry}
\usepackage{amsmath,amssymb,mathtools}
\usepackage[scr=rsfs,cal=euler]{mathalfa}
\usepackage{xfrac}
\usepackage[unicode,hidelinks,bookmarks=false]{hyperref}
\newcommand{\ie}{i.e.\ }
\newcommand{\cf}{cf.\ }
\newcommand{\resp}{respectively\ }
\newcommand{\Subsection}{\subsection}
\providecommand{\nobreakdash}{\nobreak\hbox{-}}
\setlength{\emergencystretch}{2em}
\multlinegap0pt
\usepackage{bm}
\usepackage{bbm}
\usepackage{dsfont}
\usepackage{xspace}
\usepackage{cancel}
\usepackage{mathtools}
\usepackage{amsthm}
\makeatletter
\@ifundefined{Thm}{\newtheorem{Thm}{Theorem}}{}
\makeatother
\newcommand{\aff}{ \mathbf{a} }
\newcommand{\lam}{\lambda}
\newcommand{\lest}{\leqslant}
\renewcommand{\subset}{\subseteq}
\title[Gelfand-Kirillov dimensions and annihilator varieties of hig]{Gelfand-Kirillov dimensions and annihilator varieties of highest weight modules of exceptional Lie algebras}
\author{Zhanqiang Bai, Fan Gao, Yutong Wang, Xun Xie}
\date{Source: arXiv:2509.24346v2 (CC BY 4.0)}
\begin{document}
\maketitle

\noindent\emph{Source and licence.} Gelfand-Kirillov dimensions and annihilator varieties of highest weight modules of exceptional Lie algebras. Version arXiv:2509.24346v2 (CC BY 4.0), distributed under the Creative Commons Attribution 4.0 International licence, as recorded in the arXiv licence metadata for this exact version and shown on the abstract page https://arxiv.org/abs/2509.24346v2. Full rights notice: licenses/arxiv-2509.24346-CC-BY-4.0.txt.\par\smallskip

\noindent\textit{Editorial scope.} Counted unit: the Thm whose source label is \texttt{transf} in the article (source0.tex lines 1858--1861, source label \texttt{transf}), delivered together with the article's own complete original proof of it (source0.tex lines 1863--1898, one \texttt{proof} environment of 36 lines). No mathematics is written, repaired or re-derived: statement and proof are the article's own text line for line. Required context retained uncounted: none. Every definition, display and label of those lines is retained unchanged. Omitted from this package and not redistributed: 1--1857, 1862--1862, 1899--7652. Nothing inside a retained excerpt is omitted or altered: every retained body is delivered exactly as the article's lines are written, so no omission receipt of any kind accompanies this package. Citations: the retained excerpts carry no citation command at all, so no bibliography entry of the article is transcribed and no reference is redistributed. Reference adaptation: none was needed and none was made; every reference inside the delivered unit resolves inside it, so no unresolved reference or citation is produced. Printed slips kept literal: no printed word, symbol, hypothesis or number of the counted statement or of its proof was altered. Layout and class: the article's own class is \texttt{amsart} with packages including amssymb, amsmath, amscd, amsthm, stmaryrd, verbatim, mathrsfs, tikz, color, amsfonts, enumerate, esvect, tikz-cd, dynkin-diagrams; the wrapper is the lane's pinned \texttt{amsart} editorial wrapper at 10pt on A4, which restates the theorem-like environments the retained text needs and adds the article's own macro definitions that the retained text needs (\texttt{\textbackslash aff}, \texttt{\textbackslash lam}, \texttt{\textbackslash lest}, \texttt{\textbackslash subset}), with extra wrapper packages (bm, bbm, dsfont, xspace, cancel, mathtools). The delivered numbering is the unit's own. Assets: no retained span contains a figure environment, a graphics-inclusion command, a PDF-page inclusion, a table or a TikZ picture; no image file is copied into this package, the package declares no asset dependency and no image file is redistributed with it. Licence: version arXiv:2509.24346v2 of this article is distributed under the Creative Commons Attribution 4.0 International licence, as recorded in the arXiv licence metadata for this exact version and shown on the abstract page https://arxiv.org/abs/2509.24346v2. A full rights notice is delivered at licenses/arxiv-2509.24346-CC-BY-4.0.txt. Characters: every retained excerpt is pure ASCII, so the delivered engine, XeLaTeX with its default Latin Modern text font, reports no missing character.
\begin{Thm}\label{transf}
    Suppose  $\Phi$ and $  \Psi$ are two isomorphic root systems. Let $\phi: \Phi \rightarrow \Psi$ be an isomorphism. By abuse of notation, still denote by $\phi$ the associated isomorphism  from $\mathfrak{h}_{\Phi}$ to $\mathfrak{h}_{\Psi}$, and from $\mathfrak{h}_{\Phi}^*$ to $\mathfrak{h}_{\Psi}^*$. Suppose for $\lam\in \mathfrak{h}_{\Phi}^*$, we have  an isomorphism  $\phi|_{\Phi_{[\lam]}}: \Phi_{[\lam]} \rightarrow \Psi_0 $, where $\Psi_0$ is a root subsystem of $\Psi$ and $\mathrm{rank} {(\Psi_0)}=\mathrm{rank} {(\Psi)}=n$. Then $\lam'=\phi(\lam)$ is an integral weight of type $\Psi_0$ and  we have: $$\aff(w_{\lam})=\aff(w_{\lam'}).$$

\end{Thm}

\begin{proof}
Note that if $\phi: \Phi_1 \rightarrow \Psi_1$  is an isomorphism between two root systems, then their Weyl groups $W_{\Phi_1}$  and $W_{\Psi_1}$ are isomorphic (we still denote this isomorphism by $\phi$) and we have $\phi s_{\alpha}\phi ^{-1}=s_{\phi(\alpha)}$ for any root $\alpha \in \Phi_1$  and simple reflection $s_{\alpha}$. For any $w\in W_{\Phi_1}$, we will have 
$\aff(w)=\aff(\phi(w))$. 

Recall that the Cartan matrix of a root system $\Phi$ (its simple system is $\Delta=\{\alpha_i \mid  1\lest i\lest n\})$ is defined by
$A=(A_{ij})_{n\times n}$ with $A_{ij}={\langle \alpha_i, \alpha_j^\vee \rangle}$ and $A_{ij}\in \{0,-1,-2,-3\}$ for $i \neq j$. Then we have $\alpha_j=\sum\limits_{k=1}^n A_{jk}\omega_k$. 
Suppose the simple system of $\Phi_{[\lam]}$ is $\Delta_{[\lam]}=\{\beta_1,\cdots,\beta_n\}$. For $\lam \in \mathfrak{h}_{\Phi}^*$, we can write $$\lam=\sum_{1\lest i\lest n}x_i\beta_i=\sum_{1\lest i\lest n}x'_i\omega_{\beta_i}\in \mathfrak{h}_{\Phi_{[\lam]}}^*,$$ since $\mathrm{rank} {(\Psi_0)}=\mathrm{rank} {(\Psi)}$ and $$\mathrm{span}_{\mathbb{R}}\{\omega_{\beta_1},\cdots,\omega_{\beta_n}\}=\mathrm{span}_{\mathbb{R}}\{\beta_1,\cdots,\beta_n\}.$$
Since $\lam$ is an integral weight in $\mathfrak{h}_{\Phi_{[\lam]}}^*$, we must have $x'_i\in \mathbb{Z}$ and thus
$$\lam'=\phi(\lam)=\sum_{1\lest i\lest n}x_i\phi(\beta_i)=\sum_{1\lest i\lest n}x'_i\omega_{\phi({\beta_i})}$$ is an integral weight in $\mathfrak{h}_{\Psi_{0}}^*$.


% Thus, we must have $\mathfrak{h}_{\Phi_{[\lam]}}^*=\mathfrak{h}_{\Phi}^*$ and $\mathfrak{h}_{\Psi_0}^*=\mathfrak{h}_{\Psi}^*$. 
% In other words,
% when $\lam\in \mathfrak{h}_{\Phi}^*$,  $\lam'=F(\lam)\in \mathfrak{h}_{\Psi}^*=\mathfrak{h}_{\Psi_0}^*$ will be  a weight of type $\Psi_0$.



% Then we have the following lemma.
% \begin{lemma}\label{rho}
% Let $\Pi=\{\alpha_i \mid  1\lest i\lest n\}$ be the simple roots of a finite-dimensional complex simple Lie algebra $\mathfrak{g}$. 
%    Then we have $\rho=\sum\limits_{k=1}^n \omega_k$ and 
%    $\alpha_j=\sum\limits_{k=1}^n A_{jk}\omega_k$.
% \end{lemma}

% Let $K=(k_{ij})_{n\times n}$ be the inverse of the Cartan matrix $A$. Then from Lemma \ref{rho}, we have $\omega_j=\sum\limits_{i=1}^n k_{ji}\alpha_i$. The explicit formulas for the inverses of the Cartan matrices of the simple Lie algebras are given in \cite{WZ}.





Now  we write $\lambda=w_{\lambda}\mu$, where $\mu$ is antidominant and $w_{\lambda}\in W_{[\mu]}^J$. Then $\lam'=\phi(\lam)=\phi(w_{\lambda}\mu)=\phi(w_{\lambda})(\mu)=\phi w_{\lambda}\phi ^{-1}\phi(\mu)$. From $\langle \lam, \alpha\rangle=\langle \phi(\lam),\phi(\alpha)\rangle$, we can see that $\phi(\mu)$ is still antidominant with respect to $\Psi$. If we write $w_{\lam}=s_{\beta_{i_1}}s_{\beta_{i_2}}\cdots s_{\beta_{i_k}}$ for some simple roots $\{\beta_{i_j}\mid 1\lest j\lest k\}\subset \Phi_{[\lam]}^+$, then \begin{align*}
   \phi(w_{\lam})&= \phi w_{\lam}\phi^{-1}=(\phi s_{\beta_{i_1}}\phi ^{-1})(\phi s_{\beta_{i_2}}\phi ^{-1})\cdots (\phi s_{\beta_{i_k}}\phi ^{-1})\\
    &=s_{\phi(\beta_{i_1})}s_{\phi(\beta_{i_2})}\cdots s_{\phi(\beta_{i_k})}.
\end{align*}
Thus $\aff(w_{\lam})=\aff(\phi(w_{\lam}))=\aff(w_{\lam'})$ since $\phi(w_{\lam})=w_{\lam'}\in W_{[\phi(\mu)]}^{\phi(J)}$.
\end{proof}

\end{document}
